%\newpage
%\thispagestyle{empty}
%\mbox{}
%\newpage

%%%%%%%%%%%%%%%%%%%%%%%%%%%%%%%%%%%%%%%%%%%%%%%%%%%%%%%%%%%%

\chapter{Résumé long en Français}
\chaptermark{Résumé long en Français}
\label{chap:resume_francais}

% !! utiliser la conclusion

%%%%%%%%%%%%%%%%%%%%%%%%%%%%%%%%%%%%%%%%%%%%%%
\section*{Contexte scientifique}
%%%%%%%%%%%%%%%%%%%%%%%%%%%%%%%%%%%%%%%%%%%%%%
Au sein du climat terrestre, les régions polaires sont connues pour être à la fois extrêmes et particulièrement sensibles.  
L’Arctique, constitué d’un océan central partiellement recouvert de glace de mer et entouré de terres de pergélisol, est l’une des zones les plus vulnérables aux transformations induites par le réchauffement climatique.  
Situé dans l’hémisphère Nord, il partage cette partie du globe avec les deux tiers des terres émergées, où vivent près de \perc{90} de la population mondiale.  
Cette proximité avec les zones d’activité humaine confère à l’Arctique une importance géopolitique singulière.  
Avec un rythme de réchauffement environ quatre fois supérieur à la moyenne planétaire \citep{rantanen_arctic_2022}, le climat arctique constitue à la fois un système complexe et un objet d’étude particulièrement stimulant.  
Ses conditions uniques d’exposition solaire, de stabilité atmosphérique, de température et d’isolement (présentées en \cref{sec:arctic_glimpse}) rendent son observation et sa représentation numérique particulièrement difficiles.

Le projet de recherche présenté dans cette thèse vise à combler certaines lacunes de connaissance qui limitent notre compréhension des aérosols capables de former de la glace dans les nuages, ainsi que de leurs impacts sur les nuages arctiques.  
Ces derniers, principalement des nuages stratiformes de basse altitude, exercent en moyenne un effet net de réchauffement à la surface, à l’inverse des nuages des moyennes et basses latitudes (cf. \cref{fig:matus_CRE}).  
Ils sont majoritairement composés de cristaux de glace ou d’un mélange de glace et de gouttelettes d’eau surfondue (cf. \cref{fig:matus_cloud_frac}), ce qui les rend particulièrement sensibles aux noyaux glaçogènes (INP, pour \textit{Ice Nucleating Particles}).  
Comparés aux noyaux de condensation des nuages (CCN, pour \textit{Cloud Condensation Nuclei}) \--- qui sont suffisamment abondants pour ne constituer que rarement un facteur limitant à la formation nuageuse \citep{storelvmo_aerosol_2017} \--- les INP sont plus rares de plusieurs ordres de grandeur, et sont donc déterminant pour la répartition des phases dans les nuages.

Un panache d’INP peut initier la formation de cristaux de glace dans un nuage et, aidé par l’effet de Wegener-Bergeron-Findeisen (WBF), provoquer un changement radical de phase nuageuse.  
La modification associée de l’effet radiatif des nuages (CRE, pour \textit{Cloud Radiative Effect}) peut alors affecter l’ensemble du bilan radiatif atmosphérique \citep{storelvmo_aerosol_2008}.  
La mauvaise représentation des effets des aérosols dans les modèles atmosphériques est d’ailleurs l’une des principales causes de leur difficulté à reproduire fidèlement les nuages \citep{bony_clouds_2015}.  
Le manque de connaissances sur les INP arctiques --- notamment quant aux espèces impliquées et à leurs sources d’émission --- constitue un frein majeur à une meilleure représentation de ces particules dans les modèles \citep{murray_opinion_2021}.

Bien que la science de la nucléation de glace soit étudiée depuis le XVIII\textsuperscript{e} siècle, la représentation précise de la formation des cristaux de glace dans les modèles climatiques demeure un défi important.  
Les paramétrisations de la nucléation hétérogène abondent \citep[e.g.][]{fletcher_active_1969,cooper_ice_1986, demott_predicting_2010, li_predicting_2022}, mais elles sont souvent spécifiques à une région ou à une population d’aérosols donnée, ce qui rend leur généralisation difficile dans les modèles climatiques.  
Par conséquent, ces modèles peinent encore à remplir leur rôle d’outils efficaces pour la compréhension des mécanismes atmosphériques.

Dans cette thèse de modélisation, j’ai utilisé le modèle \textit{Weather Research and Forecasting} (WRF), sa version incluant la chimie atmosphérique (WRF-Chem), ainsi que le modèle lagrangien de dispersion \textit{FLEXible PARTicle} (FLEXPART), pour explorer ces questions.  
Mon projet de recherche s’est articulé autour de trois étapes principales :  
(1) mieux comprendre les sources des INP arctiques,  
(2) améliorer la modélisation des particule capable de nucléer de la glace, et étudier comment prédire avec précision les concentrations d’INP à partir de celles-ci,  
(3) enfin, examiner l’effet des INP sur les propriétés nuageuses.


\section*{Sources des noyaux glaçogènes arctiques}
% la méthode de détection
Le premier volet de cette thèse a consisté en l'utilisation de rétro-trajectoires afin d’identifier l’origine géographique des noyaux glaçogènes (INP) observés dans l’Arctique. Pour ce faire, j’ai commencé par évaluer les performances des méthodes de rétro-trajectoires, et en particulier celles de la \textit{Potential Source Contribution Function} (PSCF).  
La PSCF, ou méthode du ratio, est une technique d’analyse appliquée aux sorties des modèles lagrangiens de dispersion de particules (LPDM), permettant d’identifier les sources probables d’émission d’espèces atmosphériques observées (\cref{P1sec:ratio_method}).  
Plus avancée que les analyses basées sur des rétro-trajectoires individuelles, largement utilisées dans la communauté des sciences de l’atmosphère, elle n’avait, à notre connaissance, jamais été évaluée de manière quantifiable.  

Afin d’évaluer la PSCF, nous avons mis en place dans WRF-Chem des émissions modélisées de traceurs atmosphériques idéalisés, et interpolé leurs concentrations sous forme de séries temporelles autour de l’Arctique afin de simuler des mesures \textit{in situ}.  
Nous avons ensuite utilisé FLEXPART-WRF pour calculer les rétro-trajectoires lagrangiennes, puis appliqué la PSCF à ces mesures simulées (\cref{P1sec:experiment}).  
Les sources d’émission étant déjà connues (puisqu’elles étaient modélisées), il a été possible d’évaluer précisément les performances de la méthode d’identification des sources (\cref{P1sec:evaluation_metric}).  
Cependant, les résultats de cette évaluation ont montré que la PSCF pouvait produire des sorties erronées et trompeuses, indiquant fréquemment des origines trompeuses (\cref{P1sec:standard_method}).  

Dans le but de construire une méthode nouvelle et précise, adaptée à l’identification de l'origine d'espèces atmosphériques à courte durée de vie en Arctique --- et en particulier des INP ---, nous avons conduit des tests de sensibilité sur les paramètres de la méthode.  
Ce travail a conduit à la définition d’une \textit{méthode du ratio améliorée}, reprenant le protocole standard de la PSCF, mais enrichie de nouveaux paramètres et de traitements supplémentaires (\cref{P1sec:improved_method}).  
Plus précisément, nous avons montré l’importance : (1)~de corriger le signal de fond saisonnier dans les séries de concentrations, (2)~de définir avec soin les seuils de percentiles identifiant les fortes et faibles concentrations dans les séries temporelles, (3)~d’établir des seuils d’exclusion pour les rétro-trajectoires non significatives, et (4)~de combiner les informations relatives aux origines probables des fortes concentrations fournies par la PSCF avec celles sur les régions identifiées comme étant dépourvues de sources d'émission.  
Enfin, nous avons évalué les limites de cette méthode améliorée afin de mieux interpréter ses résultats lors de son application à des cas réels.  
\vspace{\baselineskip}

% l’analyse des INP
Une fois cette méthode validée et ses performances connues, nous avons pu entreprendre l’identification des sources des INP observés dans l’Arctique.  
Nous avons utilisé des séries de concentrations d’INP issues de trois campagnes récentes menées dans l’Arctique central à différentes saisons (MOCCHA 2018, MOSAiC 2020 et ARTofMELT 2023, \cref{sec:INP_sources}), comprenant des mesures d'INP à différentes températures d’activation.  
Les analyses ont révélé que les INP proviennent de l’intérieur même de la région arctique, impliquant des processus de transport régionaux.  
Deux zones principales ont été identifiées : les côtes des mers de Barents et de Laptev, ainsi que les côtes du Groenland, avec une contribution possible du transport depuis l’intérieur du continent asiatique.  
Ces résultats suggèrent une prédominance des particules minérales parmi les INP observés. Cependant, l’étude des INP «~chauds~» de MOSAiC (actifs à \tempc{-10} et \tempc{-15}) a mis en évidence de possibles sources océaniques liées à l’activité biologique côtière, suggérant la présence de particules organiques marines dans les populations d’INP actifs à haute température.  
L’étude parallèle de l’origine des particules biologiques primaires (PBAP) observées durant MOCCHA 2018 a conforté cette conclusion, en indiquant également des sources océaniques issues des zones libres de glace de l’océan Arctique.  
Ces résultats encourageants motivent l’utilisation future de la méthode du ratio améliorée pour l’analyse de jeux de données d’INP passés et à venir.  


\section*{Modélisation de la nucléation de la glace et émissions d’INP}
Dans le chapitre \ref{chap4}, nous avons exploité ces nouvelles connaissances sur les sources d’INP afin d’améliorer la représentation des particules d’aérosols capables d’induire la nucléation de la glace dans WRF-Chem.  
Des émissions de poussières minérales de hautes latitudes (\cref{sec:dust_emiss}), de particules organiques marines issues des embruns marins (MPOA, \cref{sec:seaspray_emiss}), ainsi que de particules biologiques primaires continentales (PBAP) provenant de spores fongiques et de bactéries (\cref{sec:pbap_emiss}) ont été ajoutées au modèle, reproduisant de manière satisfaisante les observations de particules d'aérosol du mode grossier réalisées lors de la campagne MOCCHA 2018.  
L’objectif était d’utiliser ces particules nouvellement modélisées pour reproduire les concentrations d’INP observées \textit{in situ} dans l’Arctique central au cours de la même campagne.  
Ce défi a été abordé selon deux approches de modélisation distinctes.

\medskip
\noindent
\textbf{Première approche :}  
La première approche reposait sur la conversion, à l’aide de paramétrisations de la nucléation de la glace, des particules d’aérosols modélisées en concentrations d’INP, après leur transport dans l’atmosphère depuis leurs sources, et jusqu’au point d’observation des INP.  
Huit paramétrisations différentes ont été utilisées : trois spécifiques aux poussières minérales, une aux embruns marins, deux aux bactéries, une aux spores fongiques, et une dernière, plus générale, utilisant l’ensemble des particules du mode grossier (\cref{sec:inp_param}).  
En complément des particules modélisées, les observations de la campagne MOCCHA 2018 ont également servi d’entrée aux paramétrisations pour une évaluation des performances des paramétrisations indépendante de la modélisation des émissions de particules.  

Les résultats ont montré qu’il est difficile de reproduire la variabilité temporelle des INP à partir des particules modélisées en entrée des paramétrisations (\cref{sec:aerosols_to_INP}).  
Néanmoins, l'ordre de grandeur des concentrations est correctement reproduit par la paramétrisation utilisant la totalité du mode grossier \citep{demott_predicting_2010}.  
Lorsque les concentrations observées de particules sont utilisées en entrée des paramétrisations, la variabilité temporelle est bien mieux restituée.  
Les paramétrisations de poussières appliquées aux observations du mode grossier ont tendance toutefois à surestimer les concentrations d’INP.  
Cette surestimation s’explique par l’hypothèse sous-jacente selon laquelle les aérosols du mode grossier observés sont composés à \perc{100} de poussières minérales ; hypothèse cavalière, mais nécessaire dans l'étude pour des raisons de simplification. 
Nos résultats --- en accord avec d’autres études menées durant la campagne --- montrent d'ailleurs que le mode grossier de MOCCHA n’est pas majoritairement constitué de poussières minérales (\cref{sec:first_approach}).  

En ce qui concerne les particules biologiques observées pendant la campagne (les PBAP fluorescents ou fPBAP), aucune paramétrisation de nucléation n’a permis d’activer suffisamment de particules observées pour reproduire les concentrations d’INP mesurées.  
En effet, le taux de nucléation nécessaire pour que les INP soient dominés par les PBAP devrait approcher \perc{100} à \tempc{-20}, ce qui n'a jamais été observé et semble impossible à en croire les théories de nucléation \citep{keita_new_2020}.  
D’après l’état actuel des connaissances sur la nucléation de la glace par les PBAP, ces dernières ne peuvent, à elles seules, expliquer les concentrations d’INP observées durant la campagne.  
Le fait que les fPBAP soient bien corrélées aux mesures d’INP suggère deux possibilités : (1)~les fPBAP sont transportées conjointement avec d’autres types de particules capables d’induire la nucléation de la glace ; (2)~les observations de fPBAP sous-estiment leurs concentrations réelles.

\medskip
\noindent
\textbf{Seconde approche :}  
La seconde approche, plus novatrice, consistait à émettre directement des traceurs d’INP à partir des émissions d’aérosols considérés comme potentiellement glaçogènes, puis à simuler leur vieillissement atmosphérique.  
À l’aide de paramétrisations de nucléation de la glace adaptées à chaque type de particule (poussières minérales, PBAP, MPOA), le flux d’émission d’aérosols sont convertis en un flux d’émission d’INP réparti sur cinq classes de température, de \tempc{-5} à \tempc{-25}.  
Les traceurs d’INP sont ensuite été transportés dans l’atmosphère par le modèle en étant soumis aux processus de dépôt sec et humide, ainsi qu’à leur élimination par activation dans les nuages (\cref{sec:method_inp_emission}).  
Ce dernier processus constituait le mécanisme manquant principal dans la première approche, nécessaire pour reproduire fidèlement l’historique de nucléation des INP et le fait que les INP les plus efficaces à haute température sont préférentiellement éliminés dans les nuages froids.  

Ces traceurs d’INP ont permis d’obtenir une meilleure compréhension de la distribution spatiale des INP actifs à différentes températures, ainsi qu'une reproduction plus fidèles des observations \textit{in situ}.  
Ils ont notamment permis d’identifier les zones où la nucléation homogène peut survenir, ouvrant ainsi la voie à de nouvelles implémentations du couplage entre INP et microphysique des nuages.


\section*{Adapter la microphysique à une meilleure représentation des INP}
La dernière étape du projet de recherche présenté dans cette thèse a consisté à étudier les effets des noyaux glaçogènes (INP) sur les nuages arctiques, à l’aide du modèle atmosphérique régional WRF-Chem (\cref{chap5}).  
Les avancées réalisées dans le chapitre \ref{chap4} sur la modélisation des INP ont été exploitées pour modifier la microphysique du modèle, en utilisant directement les concentrations d’INP prédites dans la représentation de la congélation par immersion et de la nucléation par dépôt.  
Deux configurations de simulation, basées sur le schéma de microphysique sensible aux aérosols de \citet{thompson_study_2014}, ont été mises en place dans le but d'être comparées.  
Les effets de l’utilisation d’INP réalistes sur les propriétés des nuages ont été quantifiés par la comparaison d’une simulation de référence (dans laquelle les INP sont produits par la paramétrisation de \citet{cooper_ice_1986}, dite C86) et d’une simulation test (dans laquelle les INP proviennent des traceurs développés dans le chapitre \ref{chap4}).

\medskip
\noindent
\textbf{Comparaison des concentrations d’INP actifs.}  
Les premières analyses ont porté sur la comparaison des concentrations d’INP actifs entre les deux configurations.  
En effet, les traceurs d’INP développés dans le chapitre \ref{chap4} ne sont pas directement comparables aux concentrations d’INP produites par des paramétrisations dépendantes de la température, telles que C86, utilisée dans la simulation de référence.  
Les sorties de C86 incluent déjà la contrainte de température pour l’activation des INP, fournissant ainsi directement les concentrations d’INP actifs (i.e. prêts à former des cristaux nuageux).  
À l’inverse, les traceurs d’INP représentent des particules transportées dynamiquement depuis leur sources d'émission. Leur activation dépend des conditions thermiques et de leur nature chimique.  
Par conséquent, les concentrations d’INP actifs issues des traceurs dépendent à la fois des émissions, du transport et de la température d’activation.  
Une fois les conditions de température appliquées aux traceurs d’INP, les concentrations obtenues deviennent comparables à celles fournies par la paramétrisation C86.  

La comparaison montre que l’approche plus physique de la simulation à base de traceurs d’INP produit des concentrations d’INP actifs inférieures de plusieurs ordres de grandeur à celles issues de C86.  
La paramétrisation C86 tend à prédire des concentrations élevées au-dessus de la banquise et du Groenland, en raison des basses températures, même dans des zones éloignées de toute source de potentiels INP.  
À l’inverse, les concentrations d’INP étaient beaucoup plus faibles dans la simulation intégrant les informations sur le transport et le dépôt des particules.  
La distribution spatiale des INP y apparaissait moins homogène, et donc plus réaliste.  
Ainsi, en admettant la validité des traceurs développés dans le chapitre \ref{chap4} et la bonne représentation des températures par le modèle, les concentrations d’INP actifs issues de la simulation de traceurs constituent vraisemblablement une meilleure représentation des INP arctiques que celles de la simulation de référence.

\medskip
\noindent
\textbf{Comparaison des propriétés nuageuses.}  
Les propriétés des nuages simulées par les deux configurations ont ensuite été comparées.  
Dans la nouvelle configuration, l’utilisation de traceurs d’INP réalistes et de concentrations plus faibles conduit à une diminution importante et inattendue de la couverture nuageuse.  
Cette différence est la plus marquée au-dessus de la banquise, où la paramétrisation C86 tend vraisemblablement à surestimer les INP du fait de sa dépendance exclusive à la température (puisqu'il fait plus froid au dessus de la glace de mer qu'au dessus de l'océan).  
Le \textit{Ice Water Path} (IWP) présente une augmentation générale, mais une diminution locale au-dessus du Groenland et de ses côtes septentrionales.  
Quant au \textit{Liquid Water Path} (LWP), une forte diminution est observée sur le Groenland et dans l’ensemble de la région couverte par la glace de mer.  
Ces deux effets nécessitent une analyse plus approfondie d’un point de vue physique, compte tenu de la forte baisse des concentrations d’INP actifs, mais pourraient témoigner de la forte sensibilité des propriétés nuageuses aux concentrations d'INP.
Nous en avons conclu qu’il était probable que ces résultats reflètent certaines limites du schéma microphysique utilisé, qui ne parvient pas à reproduire fidèlement les processus des nuages de glace en conditions de très faibles concentrations d’INP.

\medskip
\noindent
\textbf{Perspectives.}  
Les prochaines étapes consisteront à confirmer ces constats et à examiner plus en détail les mécanismes internes du schéma de microphysique afin de l’améliorer, dans la perspective de développer un outil opérationnel permettant d’étudier la sensibilité des nuages arctiques aux INP.  
La simulation de cas d'études supplémentaires, pour lesquelles des données d’observation sont disponibles, permettra d’évaluer les deux configurations (référence et test), fournissant des indications précieuses sur la robustesse et la direction suivie par notre approche de modélisation.  
Enfin, ce \textit{set-up} de simulation constituera une base solide pour une étude détaillée du partage de phase dans les nuages et de ses impacts radiatifs associés, c\oe ur du projet d’autres doctorants actuellement en thèse au LATMOS.

%\vspace{\baselineskip}
%La présente thèse a dû être conclue avant que cette étude puisse être pleinement achevée ; néanmoins, elle ouvre la voie à de nombreuses perspectives de recherche prometteuses, déjà engagées au moment de la rédaction.
%

 
\section*{Perspectives de recherches futures}

Si les travaux présentés dans cette thèse ont apporté des réponses et des éclairages à plusieurs questions de recherche sur les interactions entre noyaux glaçogènes (INP) et nuages arctiques, ils ont également ouvert la voie à de nombreuses pistes pour de nouvelles recherches.  
Les perspectives issues de ce travail doctoral peuvent être classées en trois catégories : les prolongements et extensions des études présentées, les pistes d’exploration complémentaire, et enfin les ouvertures vers des questions plus générales.

\vspace{\baselineskip}
\textbf{Prolongements}\\
La méthode améliorée d’identification des sources d’émission de particules, présentée dans le chapitre \ref{chap3}, mérite d’être appliquée à un plus grand nombre de jeux de données de noyaux glaçogènes.  
Dans mon analyse des sources d’INP, je n’ai pas exploité l’ensemble de la campagne MOSAiC 2020, me concentrant uniquement sur la saison froide.  
Comprendre les sources d’INP au printemps et en été serait essentiel pour mieux caractériser la nature des INP en saison chaude.  
Cela viendrait compléter les résultats de l’analyse de la campagne MOCCHA 2018, et la comparaison avec les sources hivernales apporterait des éléments clés sur la saisonnalité des sources arctiques de noyaux glaçogènes.

Au-delà des données issues des campagnes sur brise-glaces, le développement de mesures d’INP de longue durée et avec une instrumentation homogène crée les conditions d’une meilleure détection des sources.  
Par exemple, le réseau mondial en expansion du \textit{Portable Ice Nucleation Experiment} (PINE) est particulièrement prometteur \citep{mohler_portable_2021}.  
PINE est un instrument automatisé et transportable mesurant les concentrations d’INP. Depuis quelques années, il est installé dans de nombreuses stations de recherche à travers le monde, y compris aux hautes latitudes.  
L’utilisation de mesures continues, couplées à des méthodes de rétro-trajectoires pour l’identification des sources, permettrait de recueillir des informations inédites sur les sources d’INP et leur variabilité saisonnière.  
L’application de la méthode d'identification des sources développée dans cette thèse à un réseau de mesures simultanées réalisées avec le même instrument, constituerait une première mondiale dans l’étude des sources d’INP.

Sur des échelles de temps plus longues, le projet \textit{Tara Polar Station}\footnote{Plus d’informations sur le site de la fondation Tara : \url{https://fondationtaraocean.org/goelette/tara-polar-station/}.} a été lancé cette année, avec pour objectif dix missions de dix-huit mois chacune dans l’océan Arctique central.  
Il fournira des concentrations d’INP mesurées à travers tout l’océan Arctique sur une période de vingt ans, offrant des informations uniques sur leur évolution spatiale et temporelle dans le contexte du réchauffement climatique, et verra probablement la disparition de la banquise estivale.  
Ces expéditions permettront d’obtenir des données sans précédent pour la détection des sources de noyaux glaçogènes, ainsi que sur leur évolution au cours du XXI\textsuperscript{e} siècle.

\vspace{\baselineskip}
\textbf{Explorations complémentaires}\\
Les deux études annexes présentées en \cref{sec:other_applications} ont démontré le potentiel d’adaptation et d’application de la nouvelle méthode de rapport dans divers contextes.  
L’étude sur l’origine des masses d’air nuageuses ouvre de nouvelles interrogations : est-il pertinent de catégoriser les nuages selon leur origine ?  
La phase des nuages dépend-elle de cette origine ?  
Les nuages en phase mixte proviennent-ils de schémas de transport spécifiques ?

Dans le chapitre \ref{chap4}, nous avons mis en place des émissions simulées de particules actualisées afin d’améliorer la modélisation des INP, mais nous ne les avons pas utilisées pour tester la sensibilité des INP à chaque type de particules.  
En activant ou désactivant les différentes espèces contribuant aux concentrations des traceurs d’INP (\cref{sec:second_approach}), il serait possible de quantifier la contribution de chaque source à la population totale d’INP, et d’analyser leur influence sur la répartition spatiale et sur le spectre de température d’activation.  
Cela justifie la mise en place de nouvelles expériences de simulation pouvant constituer une étude à part entière.

Une des lacunes de cette thèse se situe entre les résultats du chapitre \ref{chap3} et la configuration du chapitre \ref{chap4}.  
Dans \cref{sec:artofmelt_sources}, les côtes du Groenland ont été identifiées comme étant d’importantes sources d’INP, leur topographie et la nature de leurs sols suggérant des émissions de poussières minérales.  
Cependant, dans \cref{sec:dust_emiss}, le modèle d’émission de poussières utilisé ne représente pas correctement la plupart de ces sources groenlandaises.  
Cela découle des limites de la classification de l'usage des sols utilisée dans le modèle WRF-Chem, qui ne rend pas compte avec précision de la répartition des sols nus dans cette région en mutation rapide.  
Une prochaine étape importante dans la modélisation des INP consistera donc à améliorer la prise en compte des sources groenlandaises dans les émissions de particules du modèle.  
De manière similaire, les émissions d’aérosols organiques marins et continentaux reposaient sur des paramétrisations relativement simples ; des représentations plus récentes ou de nouvelles paramétrisations issues des observations arctiques récentes permettraient de mieux représenter ces sources locales d’INP dits « chauds ».

Comme mentionné à plusieurs reprises, l’étude du chapitre \ref{chap5} doit être poursuivie.  
Les questions ouvertes sur l’implémentation de concentrations réalistes d’INP dans un schéma microphysique « aerosol-aware », ainsi que les lacunes dans la connaissances de la représentation de l’impact des INP sur le climat régional des nuages, constituent l'objet du travail en cours initié durant ma thèse.

\vspace{\baselineskip}
\textbf{Questions générales}\\
Enfin, demeurent des questions plus larges que cette thèse n’a pas traitées, mais pour lesquelles elle a posé les bases.  
Dans le premier chapitre (\cref{sec:general_plan}), une question pivot a été formulée, autour de laquelle se sont articulées les études menées durant ma thèse :  
quelle est la relation causale éventuelle entre les émissions de noyaux glaçogènes, leur influence sur les nuages, et l’évolution du climat arctique ?

Mes travaux ont montré qu’il existe un lien entre les émissions d’aérosols, les populations d’INP dans l’Arctique et les propriétés nuageuses, avec de possibles conséquences importantes pour les nuages du centre de l’océan Arctique.  
Les avancées dans la compréhension des sources d’INP arctiques ont permis d’entrevoir leur sensibilité au changement climatique.  
Avec le réchauffement continu et la régression de la neige et de la banquise, il est probable que l’ensemble des sources locales d’INP identifiées dans cette thèse s’intensifient à l’avenir.  
Néanmoins, le signe de l’impact des INP sur le climat régional, tout comme l’évolution précise de leurs émissions, restent à déterminer.  
L’étude présentée dans le chapitre \ref{chap5} demeure un travail en cours nécessitant des recherches complémentaires avant de pouvoir tirer des conclusions robustes.

Je poursuivrai ces recherches aux côtés de mon collègue doctorant du LATMOS, Yaël Le Gars, dont le projet de thèse vise à améliorer la représentation de la partition de phase et des effets radiatifs des nuages arctiques en phase mixte.  
Nos résultats viendront également alimenter les travaux de Lucas Giboni, doctorant à l’Institut des Géosciences de l’Environnement (IGE) de Grenoble et co-encadré par le LATMOS.  
Sa thèse porte sur la compréhension du forçage radiatif des interactions aérosols-nuages dans le contexte du réchauffement du climat arctique.  
Son travail contribuera à une meilleure compréhension de la boucle de rétroaction entre les particules atmosphériques et le climat via leur influence sur les nuages :  
les interactions aérosols-nuages dans l’Arctique réchauffent-elles ou refroidissent-elles le climat régional, et quel en est l’impact sur les émissions d’aérosols ?

D’autres questions ont émergé au fil de mes recherches, bien qu’elles ne relevaient pas du périmètre de cette thèse.  
Parmi elles : quelle est la contribution relative de la production de glace secondaire (SIP) par rapport à la production primaire (PIP) dans la formation et l’évolution des nuages arctiques ?  
Nous nous sommes concentrés sur les processus de PIP, mais les SIP sont suspectées de jouer un rôle majeur dans la microphysique des nuages en phase mixte \citep{sotiropoulou_impact_2020, zhao_global_2021, zhao_primary_2022}.  

Enfin, l’étude du climat arctique conduit naturellement à s’interroger sur l’Antarctique.  
Bien que les INP y soient beaucoup plus rares dans l’hémisphère austral \citep{xu_characteristics_2023}, leurs sources, leur transport, leur saisonnalité, leurs effets sur les nuages et leur évolution dans un contexte de réchauffement et de dégel demeurent largement méconnus.  
Cette thèse fournit certaines des bases nécessaires pour aborder ces nouvelles questions de recherche.


% --- Boîte Questions et Réponses
\newpage
\begin{table}[H]
    \boxQAfr{Modélisation des noyaux glaçogènes arctiques}{

        \vspace{3mm}
        \centering
        \begin{tabular}{m{0.37\textwidth}|m{0.57\textwidth}}
            \textbf{Question initiale} & \textbf{Réponse apportée par la thèse} \\
            \midrule
            Quels outils sont nécessaires à l’identification des sources d’émission des particules atmosphériques ? & La modélisation de trajectoires lagrangienne, associée à des séries temporelles de mesures (in situ, aéroportées ou satellitaires) et analysée à l’aide de la méthode du ratio améliorée, permet l’identification fiable des sources d’émission de noyaux glaçogènes (INP), applicable également à d’autres types de composés atmosphériques (\cref{chap3}). \\
        \midrule
            Quelles sont les sources d’émission des INP arctiques ? & Les INP observés dans l’Arctique central proviennent probablement de l’intérieur même de la région arctique, avec un possible transport depuis l’Asie continentale. Deux principales régions sources ont été identifiées : les côtes russes des mers de Barents et de Laptev, ainsi que les côtes du Groenland. Les INP actifs à haute température (i.e. \tempc{-15} et au dessus), présentent des origines océaniques (\cref{chap3}). \\
        \midrule
            Quelles sont les espèces d’aérosols importantes pour la nucléation de glace dans les nuages arctiques ? & Les principales sources d’INP suggèrent une prédominance de particules minérales dans l’Arctique, mais l’étude des INP de la campagne MOSAiC actifs à plus haute température (\tempc{-10} et \tempc{-15}) a mis en évidence de possibles sources océaniques, indiquant la présence de particules organiques marines parmi les INP dits « chauds » (\cref{chap3}). \\
        \midrule
            Comment les différentes représentations de la nucléation de glace se comparent-elles dans un modèle régional ? & Les particules organiques simulées, combinées à des paramétrisations adaptées de nucléation de glace, produisent des concentrations d’INP plausibles.  
            Les concentrations mesurées de PBAP fluorescents (fPBAP), converties en INP à l’aide de paramétrisations de nucléation, donnent toutefois des niveaux trop faibles pour soutenir l’hypothèse d’INP dominés par les PBAP (\cref{chap4}). \\
        \midrule
            Est-il possible d’utiliser une paramétrisation générale pour les INP arctiques ? & Oui, mais cela conduit probablement à des résultats biaisés, en particulier avec les paramétrisations dépendantes de la température (\cref{chap5}). 
        \end{tabular}
    }
\end{table}

\begin{table}[H]
    \boxQAfr{\textit{Suite}}{

        \vspace{3mm}
        \centering
        \begin{tabular}{m{0.37\textwidth}|m{0.57\textwidth}}
            \textbf{Question initiale} & \textbf{Réponse apportée par la thèse} \\
            \midrule
            Est-il possible de représenter fidèlement les INP arctiques avec les connaissances actuelles, et comment procéder ? & L’émission de traceurs d’INP constitue une approche efficace pour reproduire des niveaux cohérents de concentrations d’INP.  
            Les concentrations de l’ordre de $10^{-1} \si{\per\liter}$ sont rares dans l’Arctique central pour des INP actifs à des températures supérieures à \tempc{-20} (\cref{chap4}).  
            Une fois intégrée dans la microphysique, cette approche prédit des concentrations plus faibles que les paramétrisations générales, notamment dans les régions les plus isolées de l’Arctique (\cref{chap5}). \\
            \midrule
            La présence d’INP dans l’Arctique est-elle importante pour la formation et les propriétés des nuages ? & Les nuages simulés se révèlent très sensibles à la manière dont les INP sont représentés.  
            Le traitement des INP dans les schémas microphysiques peut induire des impacts majeurs sur la fraction nuageuse et sur la phase des nuages (\cref{chap5}). \\
        \midrule
            Les INP ont-ils un effet significatif sur le bilan radiatif arctique via leurs interactions avec les nuages ? & Les travaux menés n’ont pas permis d’apporter une réponse définitive à cette question.  
            Cependant, l’influence suggérée des INP sur les nuages implique une modification importante du bilan radiatif nuageux à la surface (\cref{chap5}). \\
        \end{tabular}
    }
\end{table}

\vfill
\rule[\baselineskip]{0pt}{\baselineskip}

\newpage
\thispagestyle{empty}
\mbox{}
\newpage
